\section{Related Work}

The editing methods considered in this thesis fall into two broad families: parameter-preserving methods and parameter-modifying methods. This distinction is important because it reflects a deeper design choice about where updated knowledge should live after an edit is applied.

Parameter-preserving methods leave the original backbone weights unchanged and instead introduce auxiliary structures, such as external memory, routing components, or learned adaptors, to steer the model toward updated behavior. Their main advantage is controllability. Since the base parameters are not rewritten directly, recent edits can in principle be inserted, overwritten, or removed more explicitly. This makes such methods attractive for lifelong settings in which updates arrive incrementally and may need to be revised before they are permanently accepted. Their main weakness is that success now depends on access rather than only storage: the correct auxiliary knowledge must be activated reliably under different prompt forms, and the supporting structure must remain manageable as the number of edits grows.

Parameter-modifying methods directly update the model weights so that the edited fact becomes internalized in the same parameter space as the rest of the model's knowledge. Their main advantage is integration. If the update succeeds, the model can behave as though it had always stored the revised fact, without relying on a separate runtime memory or routing path. Their main weakness is that all old and new knowledge must coexist inside a shared and limited parameter space. As a result, repeated sequential updates become vulnerable to interference, locality loss, and cumulative instability, especially when several edits compete for overlapping representational regions.

\subsection{Parameter-Preserving Methods}

WISE~\cite{wang2024wise} represents the view that edited knowledge should be treated as side memory rather than immediately absorbed into the shared backbone. During inference, the model is guided by stored edited knowledge that is selected according to the current query. This makes WISE naturally attractive for lifelong editing because new facts can be maintained explicitly and, in principle, updated without repeatedly rewriting the same shared parameters. At the same time, its effectiveness depends strongly on whether the correct stored information is activated under rephrased prompts and whether unrelated prompts can avoid unnecessary interference from the same side memory.

GRACE~\cite{hartvigsen2023aging} also belongs to the parameter-preserving family, but it does so through learned key-value adaptors rather than a pure external memory bank. These adaptors provide a mechanism for accumulating corrections over time while keeping the original backbone intact. In conceptual terms, GRACE is closer to an online editing system than direct rewriting methods because it supports incremental accumulation of edits without repeatedly altering the same model weights. Its practical difficulty is similar to that of WISE: the stored correction is useful only if it remains reachable, selective, and non-destructive under later updates.

\subsection{Parameter-Modifying Methods}

Fine-tuning (FT) is the simplest parameter-modifying baseline in the thesis. It performs direct gradient-based updates on selected model weights so that the revised target answer is preferred on the edited prompt. This flexibility makes it a natural reference point, but it also exposes the weakness of naive sequential rewriting. When applied repeatedly, fine-tuning modifies the same shared parameter space without strong protection against interference, which makes it a useful collapse baseline in the lifelong setting.

ROME~\cite{meng2022locating} introduces a more targeted form of parameter modification. It assumes that a factual association can be represented as a localized key-value mapping in a feed-forward layer and rewrites that mapping through a rank-one update. This gives ROME a much more surgical character than ordinary fine-tuning and explains why it performs well on isolated edits. However, repeated sequential use still accumulates overlap because each new rewrite is ultimately applied to the same model parameters.

MEMIT~\cite{meng2022mass} extends the same intuition to a larger edit regime. Instead of solving a separate local rewrite for one fact at a time, it distributes the update across several layers and solves a joint system for the edited knowledge. This makes MEMIT particularly well suited to batch editing and helps explain why it remains one of the strongest methods under larger edit sets. Its cost is that the method is heavier, more dependent on model-specific assumptions, and less convenient to apply continuously than simpler baselines.

AlphaEdit~\cite{fang2024alphaedit} builds on MEMIT by adding stronger protection against destructive interference. Rather than only seeking a strong rewrite, it also constrains the update toward safer directions in parameter space so that old behavior is less likely to be damaged by later edits. This added protection improves stability in many settings, but it also creates a second trade-off: as more edits accumulate, the same protection can make later updates more restrictive or more fragile.

The six evaluated baselines therefore do not merely provide a method ranking. They instantiate two different responses to the lifelong editing problem. Parameter-preserving methods place more emphasis on controllability and incremental operation. Parameter-modifying methods place more emphasis on direct and durable factual integration. The empirical comparison in this thesis is designed to examine how that distinction appears in practice on a common evaluation setting.
